\uuid{qL4I}
\titre{Changements de variables/linéarité}
\auteur{hanani ; mzari}
\organisation{exo7}
\licence{CC BY-NC-SA 4.0}
\datecreate{2020-01-20}
\niveau{L1}
\module{Analyse}
\chapitre{Calcul d'intégrales}
\sousChapitre{Autre}
\difficulte{3}
\theme{Calculs d'intégrales, changements de variables, linéarité}
\exo7qcmid{lille-2.integrales.032}

\contenu{
\begin{question}
\qclassification{127.99}
\qtags{motcle=changements de variables/linéarité}
On pose $\displaystyle I=\int _0^{\pi}\frac{x\sin x}{1+\cos ^2x}\mathrm{d}x$ et $\displaystyle J=\int _0^{\pi}\frac{\sin x}{1+\cos ^2x}\mathrm{d}x$. Parmi les affirmations suivantes, cocher celles qui sont vraies :
\begin{answers}
\good{Le changement de variable $t=\cos x$ donne $\displaystyle J=\int _{-1}^1\frac{\mathrm{d}t}{1+t^2}=\frac{\pi}{2}$.}
\good{Le changement de variable $t=\pi -x$ donne $\displaystyle I=\pi J-I$.}
\bad{$\displaystyle I=\int _0^{\pi}x\, \mathrm{d}x.\int _0^{\pi}\frac{\sin x}{1+\cos ^2x}\mathrm{d}x=\frac{\pi ^2}{2}J$.}
\bad{$\displaystyle I=\frac{\pi ^3}{4}$.}
\end{answers}
\vskip2mm
\begin{explanations}
Avec $t=\cos x\Rightarrow \mathrm{d}t=-\sin x\, \mathrm{d}x$, on obtient :
$$J=\int _{-1}^1\frac{\mathrm{d}t}{1+t^2}=\Big[\arctan t\Big]_{-1}^1=\frac{\pi}{2}.$$
Avec $t=\pi -x\Rightarrow \mathrm{d}t=-\mathrm{d}x$ et puisque $\sin (\pi -t)=\sin t$, on obtient :
$$I=-\int _{\pi}^0\frac{(\pi -t)\sin t}{1+\cos ^2t}\mathrm{d}t=\int _0^{\pi}\frac{(\pi -t)\sin t}{1+\cos ^2t}\mathrm{d}t=\pi J-I.$$
On en déduit que : $\displaystyle I=\frac{\pi }{2}J=\frac{\pi ^2}{4}$.
\end{explanations}
\end{question}
}
